\section{Impact of Model Scale: Can Small Models Outperform Large Models?}
\label{app:model_scale}

\begin{table*}[h]
    \centering
    \small
    \begin{tabular}{lcccccc}
        \toprule
        \textbf{Model \& Method} & \textbf{AIME24} & \textbf{AIME25} & \textbf{MATH500} & \textbf{GSM8K} & \textbf{MMLU} & \textbf{Bio} \\
        \midrule
        Qwen3-32B (CoT) & 36.67 & 23.33 & 80.80 & \textbf{95.67} & \textbf{83.33} & \textbf{59.09} \\
        \textbf{Qwen3-8B (Ours)} & \textbf{36.67} & \textbf{36.67} & \textbf{83.80} & \textbf{95.67} & 83.00 & 52.72 \\
        \midrule
        \textit{Gap (8B vs 32B)} & \textit{0.00} & \textit{+13.34} & \textit{+3.00} & \textit{0.00} & \textit{-0.33} & \textit{-6.37} \\
        \bottomrule
    \end{tabular}
     \caption{Performance comparison between Qwen3-8B (with DynaDebate) and Qwen3-32B (with Standard CoT). \textbf{Bold} indicates the superior performance.}
    \label{tab:scale_comparison}
\end{table*}

A critical question in multi-agent research is whether structured collaboration can compensate for a lack of intrinsic parameter scale. To investigate this, we compare the performance of a smaller model using our framework against a significantly larger model using standard reasoning methods.

We conduct a comparative experiment using the Qwen3 family to control for training data and architecture differences:
\begin{itemize}
    \item \textbf{Small Model (Ours):} \textbf{Qwen3-8B} equipped with the \textsc{DynaDebate} framework.
    \item \textbf{Large Model (Baseline):} \textbf{Qwen3-32B} utilizing standard Chain-of-Thought (CoT) prompting.
\end{itemize}
Both models were evaluated across the same diverse set of benchmarks spanning mathematical reasoning, general knowledge, and hallucination mitigation.

The quantitative results are summarized in Table~\ref{tab:scale_comparison}.


The results reveal a distinct dichotomy based on task nature:
\begin{itemize}
    \item \textbf{Reasoning-Intensive Tasks:} On challenging mathematical benchmarks, the 8B model with \textbf{DynaDebate} significantly outperforms the 32B baseline. Most notably on AIME 2025, which features difficult out-of-distribution problems, our method achieves a \textbf{13.34\%} improvement over the 32B model (36.67\% vs. 23.33\%). Similarly, on MATH500, the smaller model surpasses the larger one by 3.0\%. This suggests that for logic-heavy tasks, the quality of the reasoning process (enhanced by debate and verification) is more decisive than raw parameter count.

    \item \textbf{Knowledge-Intensive Tasks:} Conversely, on tasks requiring extensive world knowledge (MMLU) or factual retrieval (Biography), the 32B model retains an advantage. The Biography score, in particular, drops by 6.37\% with the smaller model. This aligns with the understanding that parametric knowledge scales with model size and cannot be fully compensated for by reasoning strategies alone.
\end{itemize}

These findings demonstrate that DynaDebate effectively activates the latent reasoning potential of smaller models. It allows an 8B parameter model to achieve state-of-the-art performance on complex reasoning tasks, rivaling or exceeding models 4$\times$ its size, offering a promising direction for deploying resource-efficient reasoning systems.
